\section{Additional theory and network-validation diagnostics}
\label{supp:s3-theory-network}

This note provides robustness checks supporting the theoretical and network-validation results in the main text. In particular, we clarify (i) why the signed order parameter cancels under symmetric branch selection, (ii) why activity must be tracked jointly with polarization under activity-coupled asymmetry, and (iii) why Binder cumulant crossings are preferred over susceptibility peaks for finite-size estimation of the transition point.

\begin{figure}[H]
  \centering
  \includegraphics[width=0.92\linewidth]{s2_signed_vs_abs_q.pdf}
  \caption{Signed vs.\ magnitude order parameters under symmetric branch selection. In the symmetric regime, independent realizations select either the $+q^\ast$ or $-q^\ast$ branch with (approximately) equal probability, so the signed mean $\langle Q\rangle$ cancels even when a bifurcation exists. Using $|Q|$ (or aligning signs across realizations before averaging) recovers the expected transition curve used in the main-text figures.}
  \label{fig:supp:signed-vs-abs-q}
\end{figure}

\begin{figure}[H]
  \centering
  \includegraphics[width=0.92\linewidth]{s2_tau_ratio.pdf}
  \caption{Relaxation-time ratio between activity and polarization. We estimate characteristic relaxation times $\tau_a$ and $\tau_q$ from post-perturbation decay in representative simulations. Under activity-coupled asymmetry, activity typically relaxes faster than polarization ($\tau_a/\tau_q<1$; baseline mean $\approx 0.70$), indicating that activity is not a passive slow variable and motivating joint tracking of $(q,a)$.}
  \label{fig:supp:tau-ratio}
\end{figure}

\begin{figure}[H]
  \centering
  \includegraphics[width=0.92\linewidth]{s2_binder_fss.pdf}
  \caption{Finite-size scaling via Binder cumulant crossings. Binder cumulant curves $U_4(r;N)$ for increasing system sizes cross near the critical region and provide a stable estimator for the effective transition point in finite-$N$ simulations.}
  \label{fig:supp:binder-fss}
\end{figure}

\begin{figure}[H]
  \centering
  \includegraphics[width=0.92\linewidth]{s2_susceptibility_fss.pdf}
  \caption{Susceptibility-based estimators can be unstable at finite size. Susceptibility peaks may shift non-monotonically with $N$ and can exhibit pseudo-peaks or peak-splitting in finite simulations, motivating the use of Binder crossings as the primary estimator in the main text.}
  \label{fig:supp:susceptibility-fss}
\end{figure}

